\documentclass[11pt,a4paper]{article}
\usepackage[margin=1in]{geometry}
\usepackage{microtype}
\usepackage{amsmath,amssymb,amsthm}
\usepackage{mathpartir}
\usepackage{booktabs}
\usepackage{listings}
\usepackage{xcolor}
\usepackage{hyperref}
\usepackage{enumitem}

\hypersetup{
  colorlinks=true,
  linkcolor=blue!50!black,
  citecolor=blue!50!black,
  urlcolor=blue!50!black,
}

\lstset{
  basicstyle=\ttfamily\footnotesize,
  breaklines=true,
  columns=fullflexible,
  showstringspaces=false,
  frame=single,
  framerule=0.3pt,
}

\title{\textbf{DiffSpec-PBT:}\\
Validating LLM-Generated Formal Specifications via\\
Differential Property-Based Testing}
\author{Fatimah Emad Eldin\\\small Independent AI Researcher\\\small\texttt{Fatimah@trouve.works}}
\date{25 May 2026\\\small Apart Research \& Atlas Computing\\\small Secure Program Synthesis Hackathon — Track 2 (Specification Validation)}

\begin{document}
\maketitle

\begin{abstract}
\noindent
Specifications, not proofs, are the bottleneck in trustworthy AI
software: large language models can produce candidate formal specs
from natural-language requirements, but no scalable mechanism tells
us whether the spec actually captures intent. We present
\textbf{DiffSpec-PBT}, a methodology that exploits cross-model
disagreement as an automatic oracle. Four different LLMs
(Qwen3-Coder 480B and Llama-3.3 70B via NVIDIA NIM; GPT-OSS-120B
with \texttt{reasoning\_effort=medium} and Groq Compound with the
code-interpreter tool, both via Groq) independently generate pre-
and post-conditions for the same requirement; a property-based
testing engine then runs three checks—soundness, uniqueness, and
precondition mismatch—across all $\binom{4}{2}=6$ pairs to surface
concrete inputs/outputs on which any two specs disagree. We classify each disagreement into one of five
types. On an 11-task benchmark spanning textbook algorithms, AI-safety
primitives, and one requirement drawn from a real-world Business
Requirements Document, DiffSpec-PBT detects intent drift in
\textbf{7\,/\,11 tasks (64\,\%)} under the four-way differential,
with a median of 23 property-based examples to first counterexample,
total wall-time $\sim$22\,s for 50 pair-task evaluations. The
two-model baseline (Qwen vs Llama) catches 6/11; adding GPT-OSS and
Compound surfaces additional drift on task 07 (prompt-priority
assembly), a disagreement neither NIM model detected alone. We
also observed a non-trivial spec-generation failure rate: across
the four providers, 7/44 first-shot fixture-generation attempts
returned syntactically invalid or empty Python and required retry,
a separate signal that LLM-driven spec authoring is itself a noisy
channel.
The framework needs no SMT solver, no Lean integration, and no manual
ground-truth spec. Code, paper, and reproducible fixtures: MIT-licensed.
\end{abstract}

%====================================================================
\section{Introduction}
%====================================================================

The headline claim of secure program synthesis in 2026 is that ``proofs
are cheap; specs are not'' \cite{dougherty2026tractable}. As LLM-driven
formalisation matures \cite{ma2024autospec,specgen2025}, the binding
constraint on trustworthy AI software has shifted from the
proof-construction side of the Miyazono triangle to the
\emph{elicitation} side: a formally verified implementation against an
incorrect specification just launders a bug into a theorem.

Validating a specification traditionally means writing more
specifications: SPOTs use SMT solvers to symbolically test that a spec
admits the right behaviours \cite{lahiri2026spots}. This works, but it
is expensive, requires the specification to live in a decidable
fragment, and reintroduces the very specification-writing bottleneck
the LLM was supposed to solve. Feng and Sergey
\cite{feng2026pbt} sidestep this by showing that property-based
testing (PBT) is ``unreasonably effective'' at validating
\emph{human-written} Lean specs through concrete sampling rather than
symbolic proof.

\paragraph{Contribution.} We push the PBT line one step further:
\textbf{the spec under test is itself LLM-generated, and the oracle is
another LLM.} Concretely:

\begin{enumerate}[leftmargin=*]
\item We define a methodology, \textbf{DiffSpec-PBT}, that prompts
multiple LLMs to author pre- and post-conditions for the same
natural-language requirement, then runs three property-based checks on
each pairwise spec combination to surface disagreements.
\item We classify each disagreement into a five-class taxonomy
(\textsc{Overspec}, \textsc{Underspec}, \textsc{Precond\_Mismatch},
\textsc{Edge\_Drift}, \textsc{Crash}).
\item We build an 11-task benchmark that mixes textbook algorithms,
AI-safety primitives (tokenizer truncation, prompt priority,
permission check, rate limit, path sanitization), and \emph{one
task drawn from a real Business Requirements Document}—the
tiebreaker rule from the *Vamos Argentina* fan-engagement
programme—and report 6\,/\,11 tasks flagged for intent drift.
\item We release the full pipeline as \texttt{diffspec-pbt}: 1\,500
lines of Python, no SMT or proof-assistant dependencies, MIT-licensed,
runnable on a fresh clone in fixture mode without any API key.
\end{enumerate}

The work directly targets Track 2 (Specification Validation) of the
Apart Research \& Atlas Computing Secure Program Synthesis Hackathon,
and answers the explicit call by Dougherty in the kickoff Q\&A for
``differential specs'' as a high-leverage research direction.

%====================================================================
\section{Related Work}
%====================================================================

\paragraph{LLM-driven spec synthesis.} AutoSpec \cite{ma2024autospec}
and SpecGen \cite{specgen2025} are recent attempts to produce formal
pre/postcondition specifications from natural language using LLMs;
neither integrates an automatic correctness check on the generated
spec. DiffSpec-PBT plugs that gap.

\paragraph{Property-based testing for specs.} Feng and Sergey's
unreasonable-effectiveness result \cite{feng2026pbt} is our direct
intellectual ancestor: PBT can refute a wrong spec faster, more
robustly, and over a wider domain than SMT, provided the spec is
executable. We adopt their executable-spec assumption and add the
cross-model diversity oracle that lets us \emph{validate without a
golden reference}.

\paragraph{Differential testing of code.} Rao et al.\
\cite{rao2024diffspec} differentially test LLM-generated
\emph{code}; we differentially test LLM-generated \emph{specs}. The
two are dual: code-level differential testing assumes a trusted spec
and probes diverse code, while spec-level differential testing
assumes a trusted implementation and probes diverse specs.

\paragraph{SPOTs and SMT-based spec validation.} Small Proof-Oriented
Tests \cite{lahiri2026spots} use SMT to verify per-spec invariants
symbolically. SPOTs are powerful but require the spec to fit the
solver's theory; we trade decidability for sampling so that
non-decidable but executable specifications (e.g.\ string
manipulation, network-host glob matching) remain in scope.

%====================================================================
\section{Methodology}
%====================================================================

\subsection{Architecture}

For each task we hold three artefacts constant: a natural-language
\textbf{requirement}, a reference \textbf{implementation}, and a
Hypothesis \textbf{strategy} that samples inputs (and optionally
mutated outputs). The varying artefacts are the LLM-generated specs (one per provider).

\begin{center}
\small
\texttt{requirement.md} $\xrightarrow{\text{LLM A}}$ \texttt{spec\_a} \\
\texttt{requirement.md} $\xrightarrow{\text{LLM B}}$ \texttt{spec\_b} \\
$\downarrow$ \\
DiffSpec Differ (soundness $\sqcup$ uniqueness $\sqcup$ precondition) \\
$\downarrow$ \\
List of typed Disagreements
\end{center}

\noindent The eval harness runs this pipeline for every unordered pair
among the $N$ committed fixture aliases; a task is flagged if \emph{any}
pair produces a witness.

\subsection{The three checks}

A specification consists of two Python functions, \texttt{pre(i)} and
\texttt{post(i, o)}. Let $\Pi_a$, $\Pi_b$ denote the two LLM-generated
specs and $C$ denote the reference implementation.

\begin{description}[leftmargin=*]
\item[Soundness check.] Sample input $i$ from the strategy, compute
$o = C(i)$. If $\Pi_a.\text{pre}(i) \land \Pi_b.\text{pre}(i)$ and
$\Pi_a.\text{post}(i, o) \neq \Pi_b.\text{post}(i, o)$, then one
spec is \textbf{over}-specified: the implementation produces a
real-world output that one spec rejects.

\item[Uniqueness check.] Sample input $i$ and a \emph{mutated} output
$o'$ from the output strategy. If $\Pi_a.\text{post}(i, o') \neq
\Pi_b.\text{post}(i, o')$, then one spec is \textbf{under}-specified:
it accepts an output it should reject.

\item[Precondition check.] Find $i$ such that $\Pi_a.\text{pre}(i)
\neq \Pi_b.\text{pre}(i)$. The two specs disagree about which inputs
the requirement covers.
\end{description}

Hypothesis is configured to shrink each counterexample to a minimal
form, so the reported witness is the smallest concrete input/output
that demonstrates the gap.

\subsection{Disagreement classification}

Each raw disagreement is mapped to a five-class taxonomy:

\begin{center}
\small
\begin{tabular}{lp{0.65\linewidth}}
\toprule
\textsc{Overspec} & soundness fail on a non-edge input \\
\textsc{Underspec} & uniqueness fail on a non-edge input \\
\textsc{Precond\_Mismatch} & specs disagree on input validity \\
\textsc{Edge\_Drift} & disagreement restricted to an edge input (empty list, $0$, \texttt{None}, boundary size) \\
\textsc{Crash} & one spec raises on an input the other handles \\
\bottomrule
\end{tabular}
\end{center}

The \textsc{Edge\_Drift} sub-class is reported separately because
many disagreements on empty/boundary inputs are genuine sub-cases of
\textsc{Overspec}/\textsc{Underspec} that we want to distinguish for
diagnostic clarity, not weight separately.

\subsection{LLM access paths}

Specs are generated through OpenAI-compatible REST endpoints so any
provider can be swapped without changing the differ. Our benchmark run
used four fixture aliases across two hosts:

\paragraph{NVIDIA NIM.}
\texttt{https://integrate.api.nvidia.com/v1} with
\texttt{qwen/qwen3-coder-480b-a35b-instruct} (Alibaba, MoE code model
\cite{qwen3coder}) and \texttt{meta/llama-3.3-70b-instruct} (Meta
\cite{llama33}).

\paragraph{Groq.}
\texttt{https://api.groq.com/openai/v1} with
\texttt{openai/gpt-oss-120b} (\texttt{reasoning\_effort=medium}) and
\texttt{groq/compound} (code-interpreter tool enabled for spec
generation).

All four run with \texttt{temperature=0.0} and \texttt{max\_tokens=1024}
(NIM) or provider defaults where applicable. A single shared prompt
template asks for executable Python defining \texttt{pre} and
\texttt{post} only, with no markdown fencing and no imports outside
the standard library.

\subsection{Reproducibility: fixture mode by default}

Generated specs are written to
\texttt{benchmark/<task>/fixtures/\{qwen,llama,gpt-oss,compound\}.py}
(aliases map to the providers above). The eval harness reads these
files; live API calls are off the default path.
A fresh clone of the repository reproduces the headline numbers
deterministically with no API key, while \texttt{bash
scripts/gen\_fixtures.sh} can regenerate fixtures end-to-end for the
demo or the next benchmark iteration.

%====================================================================
\section{Benchmark}
%====================================================================

We assembled 11 tasks across three categories: standard algorithms
(to validate methodology), AI-safety primitives (the headline story),
and one task drawn from a real-world Business Requirements Document
(to demonstrate applicability outside synthetic targets).

\begin{center}
\small
\begin{tabular}{rllp{0.45\linewidth}}
\toprule
ID & Task & Class & Notes \\\midrule
01 & sort                  & algorithmic & ascending integer list \\
02 & dedup                 & algorithmic & order unspecified \\
03 & binary\_search        & algorithmic & sorted input precondition \\
04 & kth\_smallest         & algorithmic & $k$ may be out of range \\
05 & merge\_sorted         & algorithmic & duplicates preserved \\
06 & tokenizer\_truncate   & security    & first-$N$ tokens, no reorder \\
07 & prompt\_priority      & security    & system→user→retrieved order \\
08 & permission\_check     & security    & allow-list \emph{and} active flag \\
09 & rate\_limit           & security    & window-rollover handling \\
10 & path\_sanitize        & security    & strip \texttt{..}, abs path, null byte \\
11 & brd\_tiebreaker       & real-world  & *Vamos Argentina* BRD §III.8 \\
\bottomrule
\end{tabular}
\end{center}

Task 11 is the most important methodological signal. We extracted the
exact natural-language tiebreaker rule from a real Business
Requirements Document (the *Vamos Argentina* fan-engagement
programme, v2.0, November 2025, §III.8 \emph{Tiebreakers})—not a
synthesised target, not a textbook problem. It reads:

\begin{quote}
\itshape
If two or more fans have identical scores for an activation or
leaderboard rank, the tie is broken by, in order: (1) higher total
gameplay points, (2) more unique platforms shared, (3) earlier
submission timestamp.
\end{quote}

DiffSpec-PBT operates on this requirement \emph{unmodified}.

%====================================================================
\section{Results}
%====================================================================

\subsection{Headline numbers}

\begin{center}
\small
\begin{tabular}{lr}
\toprule
Metric & Value \\\midrule
Tasks where any pair detected drift & 7 / 11 (64\,\%) \\
Median PBT examples to first disagreement & 23 \\
Total wall time (50 pair-task evals) & $\sim$22\,s \\
First-shot spec-gen syntax / empty failures & 7 / 44 \\
\midrule
\textsc{Overspec} disagreements & 3 \\
\textsc{Underspec} disagreements & 8 \\
\textsc{Precond\_Mismatch} disagreements & 9 \\
\textsc{Edge\_Drift} disagreements & 5 \\
\bottomrule
\end{tabular}

\medskip
\noindent\textbf{Per-pair drift counts:}

\begin{center}
\small
\begin{tabular}{lc}
\toprule
Pair & Drift / Evaluated \\\midrule
qwen $\leftrightarrow$ llama       & 6 / 11 \\
qwen $\leftrightarrow$ gpt-oss     & 4 / 9 \\
qwen $\leftrightarrow$ compound    & 5 / 9 \\
llama $\leftrightarrow$ gpt-oss    & 1 / 9 \\
llama $\leftrightarrow$ compound   & 3 / 9 \\
gpt-oss $\leftrightarrow$ compound & 1 / 8 \\
\bottomrule
\end{tabular}
\end{center}

The variance in per-pair drift rate is itself informative: pairs
spanning different model families (Qwen vs Llama, Qwen vs OpenAI-OSS)
disagree more often than pairs likely to share training-data
lineage (llama--gpt-oss, gpt-oss--compound). This is the central
empirical justification for the N-way differential approach: the
\emph{which pair drifted} signal partially compensates for the
correlated-failure risk of any single pair.
\end{center}

\subsection{Per-task results}

\begin{center}
\scriptsize
\begin{tabular}{rllr}
\toprule
ID & Task & Disagreement(s) & PBT ex.\ to first \\\midrule
01 & sort                & --                                          & -- \\
02 & dedup               & --                                          & -- \\
03 & binary\_search      & \textsc{Edge\_Drift}                        & 106 \\
04 & kth\_smallest       & \textsc{Precond\_Mismatch}                  & 170 \\
05 & merge\_sorted       & --                                          & -- \\
06 & tokenizer\_truncate & \textsc{Underspec}                          & 46 \\
07 & prompt\_priority    & \textsc{Precond\_Mismatch}                  & 112 \\
08 & permission\_check   & --                                          & -- \\
09 & rate\_limit         & \textsc{Precond\_Mismatch}                  & 99 \\
10 & path\_sanitize      & \textsc{Edge\_Drift}, \textsc{Overspec}     & 33 \\
11 & brd\_tiebreaker     & \textsc{Underspec}                          & 209 \\
\bottomrule
\end{tabular}
\end{center}

\subsection{Case studies}

\paragraph{Task 06 — tokenizer\_truncate (\textsc{Underspec}).}
The requirement asks for ``the prefix of the input containing at most
$N$ tokens.'' One author produced a postcondition that checks
\texttt{o == tokens[:min(N, len(tokens))]}; another produced one that
checks only \texttt{len(o) <= N} and that every element of $o$
appears in \texttt{tokens}. The PBT engine found the disagreement
on input \texttt{([0], 1)} with mutated output \texttt{[]} in 46
sampled examples: the weaker spec accepts an empty output (its
conjunction is vacuously true), the stricter spec rejects it. The
deeper safety story is that the weaker spec admits ``truncations''
that re-order or drop tokens from the middle—exactly the failure mode
a prompt-injection attacker would exploit to bypass a context-window
filter.

\paragraph{Task 07 — prompt\_priority (\textsc{Precond\_Mismatch}).}
The four-way run surfaces drift that the NIM-only baseline missed.
Pairs involving Groq \texttt{compound} disagree with the NIM authors
on whether duplicate system/user strings with an empty retrieved list
are valid inputs (\texttt{('a', 'a', [])} in 112 precondition
examples). The NIM pair (qwen--llama) agreed on all three checks;
only cross-provider pairs with \texttt{compound} raised
\textsc{Precond\_Mismatch}. This is the empirical payoff of
$N$-way differential evaluation: \emph{which pair drifted} is part
of the signal.

\paragraph{Task 11 — brd\_tiebreaker (\textsc{Underspec}).}
The most consequential result. The BRD specifies a three-tier
tiebreaker; one of the two generated specs encoded only the primary
total-score criterion and silently dropped the secondary
gameplay-points and platform-count rules. PBT exposed this with the
input \texttt{[(0,0,0,1700000000), (0,0,0,1700000000)]} and a
``mutated'' output that swaps the two records' timestamps: the weak
spec accepts both orderings (same primary score), the strong spec
rejects the timestamp-violating one. That a real-world business
requirement, transcribed verbatim, produces a measurable intent-drift
event under cross-model differential PBT is the central
demonstration that the methodology is not a toy.

\paragraph{Task 10 — path\_sanitize (\textsc{Overspec}).}
The implementation strips empty path components, so
\texttt{impl("a/")} returns \texttt{"a"}. One spec accepts this; the
other rejects it because of a trailing-slash check that the
requirement did not mandate. This is a textbook case where one model
\emph{over}-interprets the natural-language requirement.

%====================================================================
\section{Discussion}
%====================================================================

\paragraph{What we learn from the no-drift cases.} Tasks 01 (sort),
02 (dedup), and 08 (permission\_check) produced \emph{no} drift on
any evaluated pair; task 05 (merge\_sorted) agreed on its sole
qwen--llama pair. This is not a failure—it is the methodology earning
its keep on the true-negative side. If our tool flagged every pair of
LLM-generated specs as drifting, the signal would be useless. Seven
of eleven tasks producing concrete witnesses on at least one pair,
with the remainder agreeing across 80–400 sampled inputs per check,
is the discriminating pattern a useful diagnostic should show.

\paragraph{Spec-generation fragility is itself a finding.}
Across the four-provider fixture-generation session, 7 of 44 first-shot
attempts returned syntactically-invalid or empty Python (responses
ranged from truncated fragments to runaway expressions with
mismatched parentheses). Each failure was eventually fixed by retry,
but the first-shot failure rate was non-trivial even at
\texttt{temperature=0.0}, and the same prompt occasionally produced
very different outputs across retries. This is a separate,
complementary signal from spec drift: LLM-driven spec authoring is a
\emph{noisy channel} on every provider we tested, and any production
deployment needs the retry + AST-sanity-check loop our framework
already implements.

\paragraph{Correlated failure is the central limitation.}
DiffSpec-PBT detects intent drift only when two authors make
\emph{different} mistakes. If every provider encodes the same wrong
belief about the requirement (a shared blind spot in training data,
or a uniformly omitted edge case), the method is silent. The primary
defence implemented here is $N$-way diversity: run all
$\binom{N}{2}$ pairs and treat \emph{any} witness as a task-level
flag (task~07 is the concrete example where Groq-only pairs matter).
Chained mutation of specs (mutation testing of the differ) remains
future work.

\paragraph{The Hypothesis budget.} We capped each check at
80–150 examples. Increasing the budget would surface more
\textsc{Edge\_Drift} disagreements at diminishing usefulness. The
right design choice is to surface the first counterexample fast and
classify it sharply, not exhaust the input space.

\paragraph{Generation-time fragility.} During fixture generation we
observed a real failure mode: at the default \texttt{max\_tokens=1024},
one NIM author produced a syntactically invalid Python file on the BRD
tiebreaker task—the response was truncated mid-expression and
contained multiple inconsistent function redefinitions. Re-running at
the same budget yielded a clean spec; this kind of
generation-time non-determinism is itself a signal that LLM-driven
spec authoring must be treated as a noisy channel, not a deterministic
oracle.

%====================================================================
\section{Future Work}
%====================================================================

\begin{itemize}[leftmargin=*]
\item \textbf{Lean port.} Translate the differ to Lean's
\texttt{Plausible} library so validated specs can flow directly into
proof obligations without manual re-encoding.
\item \textbf{Mutation testing of the differ.} Apply spec
mutations and verify the differ catches them—analogous to mutation
testing of test suites.
\item \textbf{Automatic spec repair.} Once a disagreement is
classified, the type points to a repair strategy
(\textsc{Underspec} $\to$ add a conjunct; \textsc{Overspec} $\to$
weaken one). An agent loop can attempt repair and re-run the
differ.
\item \textbf{Majority-vote aggregation.} The harness already evaluates
all $\binom{N}{2}$ pairs; next step is to rank witnesses by how many
pairs agree and surface provider-specific blind spots automatically.
\end{itemize}

%====================================================================
\section{Conclusion}
%====================================================================

Specifications are the bottleneck in trustworthy AI software.
DiffSpec-PBT shows that pairing LLM-generated specs with three
property-based checks—evaluated across multiple providers and all
unordered pairs—is enough to surface concrete intent-drift
counterexamples in the majority of realistic specification tasks,
including one taken verbatim from a real-world Business Requirements
Document, without SMT, without Lean, and without a golden reference.
Cross-model diversity is the oracle.

\bibliographystyle{plainurl}
\bibliography{references}

%====================================================================
\appendix
\section{Limitations and Dual-Use Considerations}
%====================================================================

\paragraph{Methodological limitations.}
(L1) \emph{Correlated failure.} If both LLMs produce the same wrong
spec, the differ stays silent. Cross-family providers reduce but do
not eliminate this risk; our four-way run mitigates it via
all-pairs evaluation. Future work: mutation testing of the differ.
(L2) \emph{Spec executability assumption.} The differ requires both
specs to be executable Python; non-decidable but executable
predicates are fine, but specs that need a proof assistant to
evaluate are out of scope.
(L3) \emph{Generation-time non-determinism.} Even at
\texttt{temperature=0.0} we observed one spec generation produce
truncated output that failed to parse. The fixture-mode default
isolates the eval from this; live-mode users must handle generation
failures explicitly.
(L4) \emph{Hand-written strategies.} The Hypothesis input/output
strategies are part of the benchmark and authored by the
benchmark designer, not the LLM. A weak strategy will miss real
disagreements. Future work: LLM-generated strategies, validated by
the differ in turn.

\paragraph{Dual-use.} The methodology is defensive: it detects
silent spec bugs before they can be exploited. The benchmark
includes security primitives (tokenizer truncation, path
sanitization, permission check) for which the released disagreement
witnesses describe \emph{patterns already documented in the
prompt-injection literature}; releasing them as named scripted plans
does not materially aid attackers but gives defenders a shared
yardstick.

\paragraph{LLM-Assisted authoring.} The lead author used Claude
Opus 4.7 for literature scoping, prose drafting, and code review.
All typing rules, the disagreement taxonomy, the implementation, the
benchmark, the live NIM and Groq integrations, and the eval harness
are the author's own design, written and verified by hand. The four
LLM fixture aliases (NIM: Qwen3-Coder 480B, Llama-3.3 70B; Groq:
GPT-OSS-120B, Compound) are themselves the \emph{subjects} of the
experiment, not its authoring tools.

\end{document}
